%==============================================================================
%  Otodecks - Implementation Report
%  Compiles with pdfLaTeX on Overleaf. No non-standard packages.
%
%  Figures live in figures/. A figure that is not present renders as a labelled
%  placeholder box so the document still compiles.
%==============================================================================
\documentclass[11pt,a4paper]{article}

\usepackage[T1]{fontenc}
\usepackage[utf8]{inputenc}
\usepackage{amsmath}
\usepackage[a4paper,margin=2.5cm]{geometry}
\usepackage{graphicx}
\usepackage{xcolor}
\usepackage{booktabs}
\usepackage{float}
\usepackage{enumitem}
\usepackage{listings}
\usepackage[font=small,labelfont=bf]{caption}
\usepackage[hidelinks]{hyperref}

\graphicspath{{figures/}}

\definecolor{placeholdergrey}{gray}{0.45}
\definecolor{codebg}{gray}{0.96}
\definecolor{codekw}{RGB}{0,90,160}
\definecolor{codecmt}{RGB}{100,120,100}
\definecolor{codestr}{RGB}{160,60,40}

% Long \code{} identifiers cannot hyphenate, which pushes lines into the margin.
% Let TeX loosen the paragraph slightly instead of breaking identifiers with a
% hyphen that a reader could mistake for part of the name.
\emergencystretch=3em

\newcommand{\code}[1]{\texttt{#1}}

%------------------------------------------------------------------------------
% \smartfig[width]{file stem}{caption}{label}
% Includes figures/<stem>.png if it exists, otherwise a placeholder box naming
% the file that is still missing.
%------------------------------------------------------------------------------
\newcommand{\figplaceholder}[1]{%
  \fbox{\parbox[c][3.2cm][c]{0.88\textwidth}{%
    \centering
    \color{placeholdergrey}
    \ttfamily figures/#1\\[0.6em]
    \normalfont\itshape figure not found
  }}%
}
\newcommand{\smartfig}[4][\linewidth]{%
  \begin{figure}[H]
    \centering
    \IfFileExists{figures/#2.png}%
      {\includegraphics[width=#1]{#2.png}}%
      {\figplaceholder{#2.png}}
    \caption{#3}\label{#4}
  \end{figure}%
}

\lstdefinestyle{cpp}{
  language=C++,
  basicstyle=\ttfamily\footnotesize,
  keywordstyle=\color{codekw}\bfseries,
  commentstyle=\color{codecmt}\itshape,
  stringstyle=\color{codestr},
  backgroundcolor=\color{codebg},
  numbers=none,
  frame=single,
  rulecolor=\color{codebg},
  breaklines=true,
  columns=fullflexible,
  showstringspaces=false,
  xleftmargin=0.5em,
  framexleftmargin=0.5em,
}
\lstset{style=cpp}

\hypersetup{
  pdftitle={Otodecks: A Two-Deck DJ Application},
  pdfauthor={Balasubramanian Manish},
}

\begin{document}

\begin{titlepage}
    \centering

    \vspace*{\fill}

    {\Huge\bfseries Otodecks: A Two-Deck DJ Application\par}

    \vspace{1cm}

    {\Large Object Oriented Programming\par}
    {\large Implementation Report\par}

    \vspace{2cm}

    {\large
    Balasubramanian Manish\\
    University of London\\
    \today
    }

    \vspace*{\fill}

\end{titlepage}

% \begin{abstract}
% \noindent
% Otodecks is a two-deck DJ application written in C++ with the JUCE framework. It
% plays two tracks at once with independent volume and speed, keeps a searchable
% music library that survives between sessions, provides eight hot cue slots per
% track, records either deck in isolation or the microphone to a \code{.wav} file,
% manages a panel of loop samples, and estimates the tempo of a track from its
% audio. This report covers how each requirement is implemented, why the code is
% shaped the way it is, the bugs found and fixed along the way, and what still does
% not work.
% \end{abstract}

% \vfill
% \noindent\rule{\linewidth}{0.4pt}\\[0.4em]
% \small
% \textbf{Note on figures.} Code figures in this report are screenshots taken with
% the \emph{CodeSnap} extension for Visual Studio Code. Interface figures are
% screenshots of the running application. Every fragment shown is the code as it
% stands, not a tidied-up retelling of it.
% \normalsize

\clearpage

%==============================================================================
\section{What was built, and what was already there}
%==============================================================================

The starting point was the module template: \code{Main.cpp},
\code{MainComponent}, \code{DeckGUI}, \code{DJAudioPlayer} and
\code{WaveformDisplay}. Everything else is mine.

To make that boundary checkable, every region I wrote is wrapped in \code{// My
code starts here} and \code{// My code ends here}. In the five template files
those markers surround only what I added or changed. In the nine new classes they
surround the whole file.

\begin{table}[H]
\centering
\small
\begin{tabular}{ll}
\toprule
\textbf{Template, extended by me} & \textbf{Written from scratch} \\
\midrule
\code{Main.cpp} (comment removal only)            & \code{Track} \\
\code{MainComponent.h/.cpp}                       & \code{PersistenceManager} \\
\code{DeckGUI.h/.cpp}                             & \code{MusicLibrary} \\
\code{DJAudioPlayer.h/.cpp}                       & \code{PlaylistComponent} \\
\code{WaveformDisplay.h/.cpp} (comments only)     & \code{HotCueComponent} \\
                                                  & \code{AudioRecorder} \\
                                                  & \code{SamplePlayer} \\
                                                  & \code{LoopSampler} \\
                                                  & \code{BPMAnalyser} \\
\bottomrule
\end{tabular}
\caption{Division of work between template code and code written for this assignment.}
\end{table}

Just under 4{,}000 lines including blank lines, or a little over 3{,}000 without
them. The largest single file is \code{LoopSampler.cpp} at 376 lines.

One structural decision shapes most of what follows. \textbf{The application does
not use a \code{MixerAudioSource}.} The template summed both decks through one,
which produces a single mixed buffer. That works fine until you need to record
deck one without deck two bleeding into the file, at which point there is nothing
left to tap. Each deck is now rendered into its own scratch buffer, offered to
the recorder, and added to the output by hand. That is what the mixer was doing
anyway, minus the ability to intercept a single source.
Section~\ref{sec:r4} covers it properly; it is the reason the recorder works at
all.

\smartfig{app-overview}{The application window. Two decks across the top, the
recorder status bar beneath them, and the music library and loop sampler sharing
the bottom strip.}{fig:overview}

%==============================================================================
\section{R1: Core deck functionality}
%==============================================================================

\subsection{R1A: Loading audio files into audio players}

Audio can reach a deck four ways: the deck's own LOAD button, a file dragged onto
the deck, the playlist's LOAD LEFT and LOAD RIGHT buttons, and the sampler's
``\code{> L}'' and ``\code{> R}'' buttons. All four funnel into a single
function, so there is exactly one place where a file becomes a loaded deck.

\smartfig{01-deckgui-loadfile}{\code{DeckGUI::loadFile}, the one path by which
audio reaches a deck.}{fig:loadfile}

The waveform display reads the file itself instead of asking the player for
samples, so both of them need the \code{URL}. Loading a file also adds it to the
library, because hot cues are stored on the library track, and a file dragged
straight onto a deck would otherwise have nowhere to put its cues.

\paragraph{A template bug fixed here.} In the template a dropped file went
straight to \code{player->loadURL(...)}, which loads the player but not the
waveform, so a dragged-in file played against a blank display. The drop handler
now calls \code{loadFile} like every other entry point, so the two cannot drift
apart again.

\subsection{R1B: Playing two or more tracks simultaneously}

Two \code{DJAudioPlayer} instances, two \code{DeckGUI}s, summed once per audio
block. The decks are fully independent: separate transports, separate gain,
separate resampling ratio, separate cue sets.

\smartfig{05-audio-block}{Each deck is rendered separately and summed by hand.}{fig:audioblock}

\subsection{R1C: Mixing by varying volume}

Each deck's VOL slider runs from 0.0 to 1.0 and drives
\code{AudioTransportSource::setGain}. Lowering one deck while raising the other is
the mix. The slider carries a text box showing its value to two decimal places,
so a level can be returned to deliberately instead of by eye.

\subsection{R1D: Speeding up and slowing down}

The SPEED slider drives \code{ResamplingAudioSource::setResamplingRatio} over a
range of $0.25\times$ to $3.0\times$.

$1.0$ sits at the \emph{middle} of the slider's travel instead of a
quarter of the way along, because normal speed is the value a DJ keeps coming
back to:

\begin{lstlisting}
speedSlider.setSkewFactorFromMidPoint(defaultSpeed);
\end{lstlisting}

The lower bound is $0.25$ rather than zero. A resampling ratio of zero does
not mean ``paused''. It stalls the resampler outright. \code{DJAudioPlayer}
rejects it independently of the slider, so the guard survives if the range is
ever widened:

\begin{lstlisting}
// A ratio of zero stalls the resampler rather than pausing it, so it is
// rejected along with negatives.
if (ratio <= 0 || ratio > 100.0)
\end{lstlisting}

\smartfig{deck-controls}{A single deck: transport, volume, speed, position,
waveform, the hot cue grid and the record button.}{fig:deck}

\subsection{Two further template bugs}

\paragraph{Sliders started at the bottom of their range.} JUCE initialises a
slider to its minimum. With VOL at $0$ and SPEED at $0$, a freshly opened deck
was silent \emph{and} stalled, which reads as a broken application rather than a
defaulted one. Both now start at their real values, set without notification
because the player already holds them.

\smartfig{02-deck-defaults}{Slider defaults, and putting normal speed at the
centre of the travel.}{fig:defaults}

\paragraph{\code{getPositionRelative()} divided by zero.} On an empty deck the
track length is $0$, so the division returned NaN. That NaN then propagated into
the waveform playhead and the position slider, where it was considerably harder
to diagnose than at its source.

\smartfig{03-position-nan-guard}{The empty-deck guard, added where the NaN is
created rather than where it surfaces.}{fig:nanguard}

\paragraph{The position slider is now a playhead as well as a control.} A 100\,ms
timer moves it with the audio, but it yields while the user is dragging it.
Without that check the timer and the user fight over the same control:

\begin{lstlisting}
if (!posSlider.isMouseButtonDown())
{
    posSlider.setValue(position, dontSendNotification);
}
\end{lstlisting}

%==============================================================================
\section{R2: The music library}
%==============================================================================

\subsection{R2A: Loading multiple files and displaying track details}

\code{PlaylistComponent} is a \code{TableListBox} with three columns: Title,
Duration and BPM. Files are added through ADD TRACKS, which allows a multiple
selection, or by dragging a selection of files onto the table.

\code{Track} is deliberately plain data: a path, a display name, a length, a BPM,
a date added, and its cue positions. It owns no audio resources, which is what
makes it safe to copy freely between the library, the table and the decks.
Duration is read once, when the track is added.

\smartfig{06-track-fromfile}{Reading a track's duration. A file the format
manager cannot read still becomes a row, so the user can see it and work out why
it will not play.}{fig:trackfromfile}

\code{MusicLibrary} is a \code{ChangeBroadcaster} and the playlist listens to it.
Nothing ever has to remember to tell the table to repaint; the table rebuilds
itself whenever the library says it has changed.

Beyond the requirement there is a search box, filtering by title and ignoring
case. The filter lives in the library and not in the table.

\smartfig{07-library-search}{Search lives in the library, which returns the
indices of matching tracks.}{fig:search}

The playlist stores those indices in \code{visibleRows}, which makes a table row
and a library index two different things. They are converted in one place,
\code{libraryIndexForRow}, which also absorbs the case where a repaint arrives
carrying a row number the library has since outgrown.

\smartfig{playlist}{The library table with titles, durations and detected BPM.}{fig:playlist}
\smartfig{playlist-search}{The same library filtered by the search box.}{fig:playlistsearch}

\subsection{R2B: Loading from the library to different decks}

LOAD LEFT and LOAD RIGHT send the selected track to deck one or deck two, and
double-clicking a row loads deck one. The playlist does not know that
\code{DeckGUI} exists. It calls a \code{std::function} supplied by its owner:

\smartfig{08-load-to-deck-callback}{\code{MainComponent} wires the playlist to the
decks. The sampler is wired the same way.}{fig:callback}

Neither panel holds a deck pointer or includes \code{DeckGUI.h}, which keeps both
of them independently testable and keeps the dependency pointing one way.

\subsection{R2C: Persistence across restarts}

The library is written to \code{library.json} in the operating system's
application data folder, and saved on every change. Every decision about
\emph{where} data lives and \emph{in what format} sits in
\code{PersistenceManager}. No other class builds a data path or parses a data
file.

Saving goes through a \code{juce::TemporaryFile}, so the real file is replaced
only once the write has fully succeeded.

\smartfig{09-persistence-savejson}{Saving through a temporary file. A crash
part-way through cannot truncate an existing library.}{fig:savejson}

This is not a hypothetical edge case. The library saves on almost every action
(adding a track, setting a cue, detecting a tempo), so a mid-write crash was
basically guaranteed to happen eventually if I did not guard for it.

On load, tracks whose file has since been moved or deleted are dropped instead of
appearing as rows that can never play.

\smartfig{10-library-prune}{Dead rows are pruned on load.}{fig:prune}

\smartfig{data-folder}{The application data folder: \code{library.json},
\code{sampler.json}, and the \code{samples} and \code{recordings}
directories.}{fig:datafolder}
\smartfig{library-json}{\code{library.json}, showing a saved BPM and cue
points.}{fig:libraryjson}

\paragraph{Evidence.} Three \code{.wav} files of known length were added and their
durations read back as 3.00, 5.50 and 8.25 seconds. An edited BPM and a cue point
both survived a fresh load. Search matched case-insensitively, and behaved
correctly on an empty term and on a term matching nothing. A track whose file was
deleted from disk between sessions was absent after reload.

%==============================================================================
\section{R3: Hot cues}
%==============================================================================

\subsection{R3A: Assigning cues while playing or paused}

Each deck has eight cue slots and a CLEAR ALL button, laid out by
\code{HotCueComponent}. Clicking an empty slot stores the current playhead
position.

\smartfig{12-hotcue-setcue}{Setting a cue. A position past the end of the file is
refused instead of being stored as something the deck can never seek
to.}{fig:setcue}

This reads \code{AudioTransportSource::getCurrentPosition()}, which is valid
whether the transport is running or stopped. Setting a cue therefore behaves
identically while playing and while paused, with no second code path and no
special case to get wrong.

\subsection{R3B: Modifying cues individually}

Each slot is independent. Clicking an empty slot stores the playhead; clicking a
set slot jumps to it; alt-clicking a set slot clears just that one, after which
it behaves as empty again and the next click stores a new position. Cue three can
be moved without touching cues one, two and four.

\smartfig{11-hotcue-buttonclicked}{One click handler covering all three
behaviours: set, jump, and alt-click to clear.}{fig:cueclick}

\paragraph{Why alt-click and not overwrite-on-click.} If clicking a set slot just
overwrote it, one mis-click mid-mix wipes a cue you spent time placing. I wanted
the destructive option to require a deliberate extra step, not to be the default
click behaviour.

Set slots turn orange and carry their position in the button label, so it is
possible to see where a cue points without jumping to it and interrupting the
mix.

\smartfig{hotcues-set}{The cue grid with several slots set. Orange slots show
their stored position in seconds.}{fig:cues}

\subsection{R3C: Clearing all cues}

A dedicated CLEAR ALL button per deck. It writes a full row of \code{Track::noCue}
sentinels through the same commit path as any other cue change, which keeps cue
writing and saving in one place instead of several.

\subsection{R3D: Persisting cues per track}

Cue positions are fields on \code{Track}, so they are saved and loaded with the
library. They belong to the track, so loading it onto the other deck brings its
cues along.

Reading is defensive, because a library file written by a different build may
carry a different number of cues than the current build expects. The cue count
was raised from four to eight during development, and existing libraries kept
loading without any special handling.

\smartfig{15-track-cue-persistence}{Cues are read defensively, so a library
written by an older build still loads.}{fig:cuepersist}

\subsection{A design decision: refer to the track by index, not by pointer}

My original plan had \code{HotCueComponent::setTrack(Track*)}. I changed it to
\code{setLoadedTrack(int libraryIndex)}.

Because \code{juce::Array} can reallocate as it grows, a \code{Track*} held by a
deck would go stale the moment another file got added to the library, which
happens all the time. Looking it up by index instead of caching a pointer means
that cannot happen at all, rather than me having to remember to invalidate it in
every place a file gets added.

\smartfig{14-hotcue-loadedtrack}{Resolving the track through the library each
time, so nothing is left holding a stale pointer.}{fig:loadedtrack}

A deck holding no library track returns \code{nullptr} here, and every cue
operation checks it first, so cue actions on an empty deck do nothing at all.

\subsection{Laying out eight slots}

Eight cue buttons and a CLEAR ALL in a single row would leave each button around
55 pixels wide: too narrow to read a position in, and too narrow to hit reliably
during a mix. The grid is two rows instead, with CLEAR ALL keeping a full-height
column of its own. The column count is derived from the number of buttons instead
of being hard-coded, so changing \code{Track::maxCuePoints} does not also require
a layout change.

\smartfig{13-hotcue-grid}{The cue grid layout, derived from the button count.}{fig:cuegrid}

\paragraph{Evidence.} Cues set at 1.5\,s and 4.0\,s landed at those positions, and
jumping moved the playhead correctly. An empty slot was a no-op. A cue past the
end of the file was refused. Clearing one slot left the others intact. Both set
and cleared states survived a restart, and a library saved by the four-slot build
loaded correctly into the eight-slot one.

%==============================================================================
\section{R4: Recording and the loop sampler}
\label{sec:r4}
%==============================================================================

This was the highest-risk part of the assignment and the part I would most want a
reader to look at closely.

\subsection{R4A: Recording from the decks and from the microphone}

\paragraph{Why the mixer had to go.} Recording deck one in isolation is impossible
once both decks have been summed into a single buffer. So \code{MainComponent}
renders each deck separately, offers it to the recorder, and sums the result by
hand.

\smartfig{04-render-deck}{Rendering one deck, offering it to the recorder, and
adding it to the output.}{fig:renderdeck}

\code{deckBuffer} is sized once, in \code{prepareToPlay}. The chunking loop exists
because an audio device can hand you a block larger than the size it promised,
and resizing a buffer at that moment would allocate memory on the audio thread,
which is exactly the kind of thing that produces an audible glitch under load.

\paragraph{Audio thread discipline.} Disk writing goes through
\code{AudioFormatWriter::ThreadedWriter} on a background \code{TimeSliceThread}.
The audio thread takes a \emph{try}-lock and gives up if it cannot get it.

\smartfig{17-recorder-capture}{The audio thread's capture path: a try-lock, a
block copy, and a counter. It never allocates, touches the file system or
blocks.}{fig:capture}

I did not want a recording to silently come out broken, so dropped blocks get
counted and exposed via \code{getDroppedBlockCount()} instead of just
vanishing.

The class documents its threading contract in the header, because a reader needs
to know which thread may call what before touching any of it.

\smartfig{30-recorder-threading-doc}{\code{AudioRecorder}'s threading contract,
stated in the header.}{fig:threadingdoc}

\paragraph{The microphone path.} \code{AudioRecorder} is registered as a second
device callback so that it can see the input channels and learn the device sample
rate. Because it is a secondary callback, whatever it leaves in the output buffer
gets \emph{added} to the device output. It therefore has to leave silence behind:

\begin{lstlisting}
// This recorder is a secondary device callback; the audio it is handed is
// added to the device output, so it must leave silence behind.
for (int i = 0; i < numOutputChannels; ++i)
{
    if (outputChannelData[i] != nullptr)
    {
        juce::FloatVectorOperations::clear(outputChannelData[i], numSamples);
    }
}
\end{lstlisting}

Omitting that clear doubled the output level of everything the application
played.

\paragraph{A bug worth describing, because it was ordering and not logic.}
The first deck recording produced a valid 104-byte WAV header and no audio at
all. The cause: \code{beginWriting} calls \code{stop()} first to end any
recording already in progress, and \code{stop()} disarms the deck. It was
clearing the arm flag that \code{startRecordingFromDeck} had just set. The fix is
to arm the deck only once the writer is open.

\smartfig{16-recorder-arming}{Arming the deck after the writer is open, not
before.}{fig:arming}

\code{stop()} decides whether it is actually ending a real recording before it
tears down the writer, which is why \code{onRecordingFinished} only fires once,
and only after the file has been flushed and closed. Nothing downstream ever sees
a half-written file.

\smartfig{18-recorder-stop}{Stopping: decide first, then tear down, then
announce.}{fig:stop}

\smartfig{recording-active}{Recording deck one. The REC button turns red, the
other deck's is disabled, and the status bar names the file being
written.}{fig:recactive}
\smartfig{recording-saved}{After stopping. The status bar shows the saved path and
reveals the file in Finder when clicked.}{fig:recsaved}

\subsection{R4B: Displaying and playing back recorded loops}

\code{LoopSampler} is the panel: a table of samples with ADD, REMOVE, PLAY, STOP,
a loop toggle, a gain slider, and ``\code{> L}'' and ``\code{> R}'' to send a
sample to either deck. Double-clicking a row previews it within the panel.
\code{SamplePlayer} handles the playback and is deliberately lighter than
\code{DJAudioPlayer}, with no speed control and no seeking, because a sample pad
needs neither.

A finished recording is taken into the sampler automatically, so a loop captured
off a deck can be triggered straight back over the mix without the user going to
find the file.

\smartfig{19-recording-to-sampler}{A finished recording reaches both the library
and the sampler panel.}{fig:rectosampler}
\smartfig{20-sampler-addrecorded}{The sampler's side of it. The recording itself
stays in the recordings folder and the sampler works from its own copy, exactly
as it does for any other added file.}{fig:addrecorded}

Retriggering a sample restarts it from the beginning.

\smartfig{22-sampleplayer-play}{Retriggering restarts from zero, so a sample
behaves like a pad instead of resuming where it left off.}{fig:sampleplay}

I mix the sampler's preview in \emph{after} the recorder taps the signal, on
purpose. Otherwise auditioning a sample while recording would bleed into the
file. I left a comment there, since it would be an easy thing to accidentally
reorder later.

\smartfig{sampler-panel}{The loop sampler. The row is a loop recorded from a
deck, which reached the panel on its own when the recording finished.}{fig:sampler}

\subsection{R4C: Persisting and managing the sample list}

The sample list is its own file, \code{sampler.json}, kept separate from the music
library. Samples are loops and one-shots, not tracks, and mixing them into the
playlist would make the library worse at the job it exists to do.

Files added through the interface, dropped onto the panel, or arriving as a
finished recording are \emph{copied} into the application's \code{samples}
folder.

\smartfig{21-sampler-copyin}{Samples are copied into the samples folder. A file
already living there is left alone, so re-adding one does not pile up
copies.}{fig:copyin}

The sample set therefore survives the original file being moved or deleted. On
startup, \code{load()} reads \code{sampler.json} and \emph{then} scans the samples
folder, so a file dropped into that folder by hand also appears.

REMOVE only takes the row out of the sample list, it does not touch the file on
disk. I did not want a table button to be the thing that quietly deletes
someone's audio. Anything the format manager cannot read is refused when it is
added, so it never becomes a row that will not play.

\paragraph{Evidence.} Two hundred blocks pushed through the real render loop
produced a readable two-channel, 44.1\,kHz WAV of exactly the expected sample
count, with zero dropped blocks. Its peak level matched the deck's own output,
which is also what proves the unarmed second deck never leaked in. The microphone path was confirmed by hand. For the sampler: play, stop
and retrigger all behaved as intended; audio continued past the end of the file
when looping and stopped when not; non-audio files and duplicates were refused;
and the list survived a restart.

%==============================================================================
\section{R5: Beats Per Minute}
%==============================================================================

\subsection{R5A: What BPM is, and how a DJ application uses it}

BPM is a single number giving a track's tempo in beats per minute. In a DJ
application it is the number that decides whether two tracks can be played
together at all.

\paragraph{Scenario one: choosing the next track.} A DJ playing a 128\,BPM house
track needs the next one to be close, because a deck can only be stretched so far
before it stops sounding right. Roughly $\pm$6\% is the usual working tolerance.
Beyond that, vocals are audibly pitched up or down and drums lose their weight.
With a BPM column in the library, the DJ scans or sorts and finds a candidate in
a couple of seconds. Without one they are auditioning tracks in headphones under
time pressure, which is precisely the moment you do not want to be guessing. In
Otodecks this is the BPM column in the playlist, and that $\pm$6\% maps directly
onto the SPEED slider.

\paragraph{Scenario two: beatmatching a transition.} To mix two tracks together
their beats have to land on top of one another. That needs the same tempo, and
alignment of the beats within it. BPM gives the first of those. If deck one is at
128 and the incoming track is 124, the DJ knows before touching anything that the
incoming deck needs a ratio of $128/124 \approx 1.032$, sets it, and is left with
only the alignment to nudge. Working that out by ear in real time takes practice
most people do not have. Doing it from a BPM number is just division. Combined
with hot cues, a
DJ can drop the incoming track from a cue placed on a downbeat at a tempo that is
already matched, which is how live remixing is done.

\smartfig{bpm-column}{Detected tempo shown in the library's BPM column. Tracks
that have not been analysed show a dash instead of a fabricated
number.}{fig:bpmcolumn}

\subsection{R5B: Alternative techniques and what commercial applications do}

BPM on its own is limited. It gives the tempo but not \emph{where} the beats fall,
and it assumes the tempo never changes. The techniques that address this:

\begin{itemize}[leftmargin=1.4em]
  \item \textbf{Beat grids.} Rekordbox, Serato and Traktor store a grid of beat
    positions across the whole track, anchored to a downbeat, as well as a BPM.
    This is strictly more information. It survives tempo drift in live-recorded or
    older material, and it lets the software align \emph{phase} and not just rate.
    It is also why those applications let the user drag the grid: automatic
    detection places the beats roughly right and a human corrects the anchor.
  \item \textbf{Sync buttons.} Built on beat grids, these match tempo and phase in
    a single press. What the DJ gives up is control. When the grid is wrong, sync
    is confidently wrong.
  \item \textbf{Key detection and harmonic mixing.} Camelot or Open Key notation.
    Two tracks at the same tempo can still clash if their keys fight, so modern
    libraries show key alongside tempo and DJs select on both.
  \item \textbf{Time-stretching instead of resampling.} Otodecks changes speed by
    resampling, which changes pitch with it, so a track sped up 6\% is also 6\%
    sharp. Commercial applications use a phase vocoder or something similar to
    change tempo while holding pitch, sold as ``Master Tempo'' or ``Key Lock''.
    This is the largest functional gap between this application's speed control
    and a professional one.
  \item \textbf{Tap tempo.} The manual fallback: the DJ taps the beat and the
    application averages the intervals. It is slower and less precise, but it
    works on anything at all, including material with no clear percussive onsets
    where automatic detection simply fails. Otodecks implements this; see
    Section~\ref{sec:r5d}.
\end{itemize}

\subsection{R5C: How BPM would be calculated, conceptually}

Tempo estimation is three problems in sequence. Find where the events are, find
the period at which they repeat, and turn that period into a number.

\paragraph{Stage one: onset detection.} Musical beats usually coincide with a
sharp rise in energy: a kick, a snare, a chord stab. Reduce the audio to a
low-rate envelope by taking the RMS energy of each short frame, then take the
difference between consecutive frames. The important part is to half-wave rectify
that difference, keeping only the \emph{rises}. A falling edge is the tail of a
hit that has already been counted, and counting it again would place a spurious
onset just after every real one.

\paragraph{Stage two: periodicity.} The onset envelope is a spiky signal that
repeats at the beat period. Autocorrelation, which correlates the signal with a
time-shifted copy of itself, peaks at the lags where those spikes line up, and
the strongest peak within a plausible tempo range gives the beat period. Two
alternatives are a Fourier transform of the onset envelope, which finds the same
periodicity in the frequency domain, and detecting discrete onset times and
taking the median gap between them.

\paragraph{Stage three: conversion, and the octave problem.} A lag in frames
becomes seconds via the hop size and sample rate, and then $\text{BPM} = 60 /
\text{period}$. The hard part is that autocorrelation cannot distinguish 70 from
140. If spikes line up at a lag of $N$ frames, they also line up at $2N$. This is
the ``octave error'', and it is the most common failure mode in tempo detection.
Fixing it needs a second, independent piece of evidence, not a better peak
search.

\subsection{R5D: Implementation and visualisation}
\label{sec:r5d}

\code{BPMAnalyser} implements the method above. Pressing \textbf{BPM} in the
playlist analyses the selected track and writes the result into the BPM column,
where it is saved with the library and survives a restart.

\paragraph{The envelope and onset strength.} RMS energy per 256-sample hop,
averaged across up to two channels. A mono sum is enough for onset detection and
costs less than treating the channels separately. The difference is then
rectified.

\smartfig{23-bpm-onset}{Onset strength: only rises in energy count.}{fig:onset}

\paragraph{Autocorrelation, bounded and normalised.} Only the lags a real tempo
could occupy are searched, and each is divided by its overlap length.

\smartfig{24-bpm-autocorrelate}{Normalised autocorrelation over the tempo
range.}{fig:autocorr}

Without that normalisation the method is systematically biased towards fast
tempos, because a short lag simply has more terms in its sum than a long one.

\paragraph{Sub-frame precision.} At a 256-sample hop, one frame is worth a couple
of BPM, and the true period rarely lands exactly on a frame boundary. Fitting a
parabola across the winning peak and its two neighbours recovers the fractional
part.

\smartfig{25-bpm-parabolic}{Parabolic interpolation across the correlation
peak.}{fig:parabolic}

\paragraph{The octave check.} A second estimate is computed independently, from
the median gap between detected onset peaks. If the two disagree by roughly a
factor of two, the interval estimate wins.

\smartfig{27-bpm-intervals}{The independent estimate, from the median gap between
onsets.}{fig:intervals}
\smartfig{26-bpm-octave-check}{The cross-check that resolves half- and
double-tempo errors.}{fig:octave}

The median is used because a single missed or doubled onset drags a mean badly
and barely moves a median.

\paragraph{Tap tempo.} The \textbf{TAP} button beside BPM implements the manual
fallback discussed in R5B. \code{BPMAnalyser::registerTap} timestamps each tap and
discards a stale count on its own, and \code{getBPMFromUserTap} averages the
intervals once three taps have arrived.

\smartfig{28-tap-tempo}{Tap tempo. The button reports its count from the first
click, and the tempo once it has one.}{fig:tap}

The button gives feedback from the very first click, because two silent clicks
read as a dead button. The result is written onto the selected track on
\emph{every} tap, so the BPM column tightens as the user keeps tapping, while the
write to \emph{disk} waits until the sequence goes quiet. Saving on every tap would mean a file write twice a second while somebody
taps.

\smartfig{29-tap-finish}{Closing off a tap sequence: save once, then reset.}{fig:tapfinish}

A \code{Timer} matched to the analyser's own 2.5\,second tap timeout is what ends
the sequence, so the button and the analyser can never disagree about whether a
count is still live.

\smartfig{tap-tempo}{The TAP button mid-count, and the resulting tempo in the BPM
column.}{fig:tapshot}

Both routes fold their result into the range 70--180\,BPM, which is why tapping a
60\,BPM ballad reports 120. A half-time and a double-time reading describe the
same tempo, and folding them makes values comparable across the library.

\paragraph{Evidence.} Against synthesised click tracks of known tempo: 90 gave
89.93, 120 gave 120.12, 128 gave 127.78 and 140 gave 139.80. A 100\,BPM track
with deliberately offbeat hats, built specifically to tempt a double-tempo error,
gave 100.14. All within 0.25\,BPM. Files that cannot be read return $0$ instead of
a guess, and the playlist shows a dash. Tap tempo at 500\,ms intervals gave
119.38\,BPM.

%==============================================================================
\section{Assumptions and decisions}
%==============================================================================

\paragraph{Data lives in the OS application data folder, not the project folder.}
My original plan put everything in a \code{data/} directory beside the executable.
That does not survive being run from a \code{.app} bundle on macOS, and it is the
wrong location on Windows. It is now
\code{\textasciitilde/Library/Application Support/Otodecks} and
\code{\%APPDATA\%\textbackslash Otodecks}, decided in one place:

\begin{lstlisting}
// On macOS that special location is ~/Library itself, not ~/Library/Application Support.
#if JUCE_MAC
root = root.getChildFile("Application Support");
#endif
\end{lstlisting}

\paragraph{The library stores paths, not copies of your music.} Otodecks never
moves or copies a track file; it stores the path. The sampler is the deliberate
exception, because a sample set that breaks when you tidy your Downloads folder is
not much of a sample set. That is also why removing a sample from the panel does
not delete the file.

\paragraph{Files loaded onto a deck join the library automatically.} The library
therefore grows as you play, and grows again when a recording finishes. The alternative was a dragged-in file
having nowhere to persist its hot cues, which seemed clearly worse.

\paragraph{Panels do not know about decks.} \code{PlaylistComponent} and
\code{LoopSampler} reach the decks only through \code{std::function} callbacks set
by \code{MainComponent}. Neither holds a deck pointer or includes
\code{DeckGUI.h}.

\paragraph{Anything reachable from the audio callback avoids allocation, locks,
logging and file access.} Values crossing the thread boundary are
\code{std::atomic} or pass through the writer's FIFO. \code{deckBuffer} and
\code{sampleBuffer} are sized in \code{prepareToPlay}, and oversized blocks are
chunked so nothing triggers a resize.

%==============================================================================
\section{How this was tested}
%==============================================================================

Each stage was verified by a temporary console target compiled against the same
sources, which drove the real classes headlessly and was deleted once it passed.
That is how the recorder's arming bug and the NaN guard were found. Both would
have been unreliable to surface by clicking around the interface, and the
recording bug in particular produced a file that looked entirely plausible in a
directory listing.

The fade tail from \code{AudioTransportSource}'s gain ramp was tripping my
silence check right after \code{stop()} and \code{setGain(0)}: three failures that
were bad assertions, not bugs. Worth remembering that a failing test is not
automatically proof the code is wrong.

The application builds under two toolchains from the same \code{Source/} tree:
CMake and CLion on macOS (clang, JUCE 8.0.4), and Projucer and Visual Studio on
Windows (MSVC, JUCE 9.0.2). That was worth the effort for a reason beyond
portability. Clang caught missing destructor \code{override}s that MSVC did not,
and MSVC caught a variable in my own code shadowing \code{Component::flags} that
clang did not.

Under JUCE's recommended warning flags, every remaining compiler warning is in
one of the five template files: a \code{double} row height passed to
\code{setBounds}, destructors without \code{override}, unused parameters, and two
deprecated JUCE calls. The nine classes written for this assignment compile
silently on both toolchains.

%==============================================================================
\section{Known limitations}
%==============================================================================

What the application still does badly:

\begin{itemize}[leftmargin=1.4em]
  \item \textbf{BPM analysis blocks the message thread.} A long track freezes the
    interface for a moment, and a wait cursor is the only warning the user gets.
    Moving the analysis to a background thread is the correct fix and the obvious
    next step.
  \item \textbf{Only the first 60 seconds of a track are analysed.} Fine for
    material at a constant tempo, wrong for anything that changes tempo part-way
    through.
  \item \textbf{Speed change is resampling, so it changes pitch with it.} There is
    no key lock, as discussed under R5B.
  \item \textbf{The library grows as you use the application}, because deck loads
    and finished recordings are both added to it.
  \item \textbf{Sample removal is list-only.} REMOVE forgets a sample but leaves
    the copied file in the samples folder, so the folder accumulates over time.
  \item \textbf{JUCE 7 or newer is required.} The recorder uses
    \code{audioDeviceIOCallbackWithContext} and the sampler uses
    \code{RangedDirectoryIterator}.
\end{itemize}

%==============================================================================
\appendix
\section{Building and running}
%==============================================================================

\paragraph{CMake (CLion or the command line).} JUCE is located automatically, in
order: \code{libs/JUCE} inside the project, then \code{../JUCE} alongside it,
otherwise fetched from GitHub at configure time. Override with
\code{-DJUCE\_PATH=/path/to/JUCE}.

\begin{lstlisting}[language=bash]
cmake -B build
cmake --build build
\end{lstlisting}

\paragraph{Projucer (Xcode or Visual Studio).} Open \code{OtoDecks.jucer}, set the
global JUCE modules path to a JUCE 8 checkout, then save and open the generated
project. Both build systems reference the same files in \code{Source/}, so they
cannot drift apart.

\paragraph{Where data is stored.}

\begin{table}[H]
\centering
\small
\begin{tabular}{ll}
\toprule
macOS   & \code{\textasciitilde/Library/Application Support/Otodecks/} \\
Windows & \code{\%APPDATA\%\textbackslash Otodecks\textbackslash} \\
\bottomrule
\end{tabular}
\end{table}

\noindent
containing \code{library.json}, \code{sampler.json}, \code{samples/} and
\code{recordings/}. Deleting these resets the application to an empty state.

\end{document}
